% Author: Izaak Neutelings (February 2020)
\documentclass[border=3pt,tikz]{standalone}
\usepackage{physics}
\usepackage{xcolor}
\usepackage{calc}
\usetikzlibrary{calc}
\usetikzlibrary{decorations.markings}
\tikzset{>=latex} % for LaTeX arrow head

\colorlet{Ecol}{orange!90!black}
\colorlet{Bcol}{violet!90}
\colorlet{Icol}{blue!70!black}
\colorlet{gausscol}{green!40!black}
\colorlet{gausscol2}{green!45!blue}
\tikzstyle{current}=[->,Icol,thick]
\colorlet{pluscol}{red!60!black}
\colorlet{minuscol}{blue!60!black}
\tikzstyle{anode}=[top color=red!20,bottom color=red!50,shading angle=20]
\tikzstyle{cathode}=[top color=blue!20,bottom color=blue!40,shading angle=20]
\tikzstyle{gauss surf}=[gausscol,top color=green!2,bottom color=green!80!black!70,shading angle=5,fill opacity=0.4]
\tikzstyle{metal}=[top color=black!15,bottom color=black!25,middle color=black!20,shading angle=10]
\tikzstyle{mydashes}=[dash pattern=on 1 off 1]
\tikzset{
  EFieldLine/.style={thick,Ecol,line cap=round,decoration={markings,
                     mark=at position #1 with {\arrow{latex}}},
                     postaction={decorate}},
  BFieldLine/.style={thick,Bcol,postaction={decorate},decoration={markings,
                     mark=at position #1 with {\arrow{latex}},
                     mark=at position #1+0.5 with {\arrow{latex}}}},
  EFieldLine/.default=0.5,
  BFieldLine/.default=0.4}
\usetikzlibrary{3d}

\begin{document}


% CAPACITOR 3D - displacement current derivation
\begin{tikzpicture}[scale=1.5]

    
    \begin{scope}
        
        \coordinate (P2) at (0, 0);
        \coordinate (P1) at (4, 2);
        


        \draw[gausscol, thick, -latex] (P2) node [below, black] {$P_2$} -- (P1) node [midway, below] {$\boldsymbol{r}$};
        \filldraw[gausscol] (P2) circle (1pt);
        \draw[minuscol, very thick, -latex] (P2) -- ++ (0, 0.3) node [above] {$\ddot{\boldsymbol{p}}$};
        \node [right] at (P1) {$P_1$};
        \draw[Bcol, very thick, -latex] (P1) -- ++ (-0.3, +0.13) node [left, yshift=3, xshift=2] {$\boldsymbol{B}$};

        \draw[Ecol, very thick, -latex] (P1) -- ++ (0, -0.4) node [below, yshift=3, xshift=2] {$\boldsymbol{E}$};


        % \draw[gausscol, thick, -latex] (P2) node [left, black] {$P'_2$}-- (R) node [midway, above] {$\boldsymbol{r}'$};

        % \node[right] at (R) {$P_1$};

        % \draw[Ecol, very thick, -stealth] (P2) -- ($(P2)!0.2!(R)$) node [above] {$\boldsymbol{e}_{r'}$};
        % \draw[Bcol, very thick, -stealth] (P2) -- ++(-48:0.5) node [below] {$\boldsymbol{v}$};
        % \filldraw[] (P2) circle (1pt) node [above] {$q$};
        % \filldraw[] (P1) circle (1pt) node [above] {$q$};


        % \draw[pluscol, thick, ->] (1, -1.7) node[fill=white] {Position at $t$} to[bend left=30] (P1);
        % \draw[pluscol, thick, ->] (-1, -1) node[fill=white] {Position at $t-r'/c$} to[bend left=30] (P2);


 
    \end{scope}

    
\end{tikzpicture}
  
\end{document}